\chapter{Methodology}\label{ch:method}
%% ---------------------------------------------------------------------
%% PURPOSE: describe your approach precisely enough that a competent
%% reader could reimplement it.
%%
%% Rules of thumb:
%%   - Start with an overview figure and one paragraph walking through it.
%%   - Then one section per component, in the order data flows through them.
%%   - For each design decision, state the alternative and why you rejected it.
%%   - Use the running example from the background chapter.
%%   - Separate WHAT the method does from HOW you implemented it
%%     (implementation details: a section at the end, or an appendix).
%%   - Pseudocode (algorithm environment) for anything with a loop.
%% ---------------------------------------------------------------------

\section{Overview}\label{sec:overview}
\todo{Overview figure + walkthrough.}

\begin{figure}[htbp]
  \centering
  \fbox{\parbox[c][4cm][c]{0.7\textwidth}{\centering Overview diagram goes here}}
  \caption{Overview of the approach. \todo{A caption should be readable without the text.}}
  \label{fig:overview}
\end{figure}

\section{\todo{Component 1}}\label{sec:component-1}
\todo{Describe the component.}

\begin{algorithm}[htbp]
  \caption{\todo{Algorithm name}}\label{alg:main}
  \begin{algorithmic}[1]
    \Require input $x$
    \Ensure output $y$
    \State $y \gets x$
    \While{not converged}
      \State refine $y$
    \EndWhile
    \State \Return $y$
  \end{algorithmic}
\end{algorithm}

\section{Implementation}\label{sec:implementation}
\todo{Languages, libraries, versions, hardware. Link to the code.}
